\documentclass[a4paper,11pt,openany]{report}
\usepackage{../preamble/agenticboost}

% Source evidence: supplied conversation and refactor transcript.
% ChatGPT file citation for the evidence packet: fileciteturn68file0
%
% Intended repository path:
% agents/reports/tex/sprints/agentic-system-refactor-retrospective.tex
%
% This report is human-facing evidence. It is not an operational instruction
% source for future agents. Current rules remain in agents/rules/core/, agents/rules/workflows/,
% and agents/rules/roles/.

\renewcommand{\sprintId}{AgenticSystem Refactor}

\newcommand{\safePath}[1]{\texttt{\detokenize{#1}}}
\newcommand{\decision}[1]{\textbf{#1}}
\newcommand{\risklabel}[1]{\textsc{#1}}

\title{%
  \textsc{AgenticCareerBoost}\\[0.4em]
  \Large AgenticSystem Refactor Retrospective\\[0.3em]
  \large Failure Analysis $\cdot$ Cleanup Corrections $\cdot$ Minimal Architecture Rules%
}
\author{Dídac Llorens \\ \small \href{https://github.com/DidacLL}{github.com/DidacLL}}
\date{July 2026 --- v0.1}

\begin{document}
\maketitle
\tableofcontents
\newpage

% ============================================================================
\chapter{Executive Summary}
\label{ch:executive}
% ============================================================================

\takeaway{The AgenticSystem refactor did not fail because the objective was
wrong. It failed because the system kept solving control problems by adding more
control surface. The useful result is a clearer architecture: rules live in
\safePath{agents/rules/}; state records evidence under \safePath{agents/state/};
drafts are candidate outputs; public reports are proof artifacts; the site owns
runtime files and downloadable public artifacts; and validation must match the
requested execution mode.}

This report documents the failure path that preceded the AgenticSystem cleanup,
the mistakes introduced during the cleanup itself, the corrections made through
human review, and the architecture that emerged from those corrections.

The purpose is not to preserve a transcript. The purpose is to turn a painful
refactor sequence into reusable engineering evidence. The project repeatedly
exposed the same pattern: agents expanded narrow work into process work, then
validated the process rather than the user's requested output. When the user
asked for content, the system drifted toward tests, state updates, site edits, or
new documentation. When the user asked for cleanup, the system initially treated
historical evidence as clutter. When the user asked for less structure, the
system proposed new folders and generated deployment layers.

The corrective principle is therefore strict:

\begin{quote}
Do not add structure to compensate for unclear authority. First define authority.
Then delete duplication. Only then change architecture.
\end{quote}

The final target is a smaller system, not a weaker one. The project should keep
the artifacts that prove work happened: formal reports, state logs, summaries,
research, and public pages. It should remove or demote sources that can steer
future agents incorrectly: stale social drafts, superseded style books, old tool
outputs, manual mirror copies, and template/benchmark rules that reward ceremony
instead of behavior.

\section{Primary Conclusions}

\begin{enumerate}
  \item \textbf{The instruction hierarchy must be operationally enforceable.}
    Saying that direct user instructions outrank repository rules is not enough
    if workflow files still trigger escalation, tests, review loops, or state
    writes by default.

  \item \textbf{State is evidence, not law.} Agents may inspect
    \safePath{agents/state/} for recent context, but they must never treat it as the
    source of behavior rules, voice rules, acceptance criteria, or future-run
    strategy.

  \item \textbf{Execution mode is the missing boundary.} A text-only task, a
    site-copy task, an implementation task, and a system-refactor task require
    different allowed writes and different validation. Without this boundary,
    ``where appropriate'' becomes a permission slip for overwork.

  \item \textbf{Human-facing evidence is not clutter.} Deleting logs, summaries,
    or LaTeX reports removes the proof body of the project. The right fix is to
    make evidence non-authoritative, not to erase it.

  \item \textbf{The site must not maintain tracked mirror copies.} If the repo
    already contains the canonical file, the site should reference it directly
    under the deployment model or expose the repository root as the static
    surface. Manual copied files are a guaranteed drift source.

  \item \textbf{Validation must test the changed behavior.} Link checks and
    Python tests do not prove the browser still works. Runtime changes require a
    browser/load check, while text-only work should not run broad technical gates
    unless explicitly requested.
\end{enumerate}

\pitfall{The most dangerous failure was not one bad deletion or one bad site
change. The dangerous pattern was the system's habit of answering architecture
questions by adding more architecture.}

% ============================================================================
\chapter{Evidence Scope and Reading Rules}
\label{ch:evidence-scope}
% ============================================================================

\takeaway{This report is an evidence artifact. It summarizes the refactor
conversation and the resulting design conclusions. It does not replace the
canonical operating rules in \safePath{agents/rules/core/}, \safePath{agents/rules/workflows/},
or \safePath{agents/rules/roles/}.}

\section{Evidence Inputs}

The report is based on the conversation evidence body created during the
AgenticSystem cleanup/refactor discussion. That evidence contains three distinct
layers:

\begin{description}[style=nextline]
  \item[Original symptoms]
    Documentation loops, copy drift, tests running for text-only work, unwanted
    site-code edits, duplicated site files, and old instruction files poisoning
    future runs.

  \item[Refactor attempts]
    Execution-mode planning, public-copy canon planning, agents/state/evidence cleanup,
    site-source deduplication, generated artifact proposals, root-upload
    alternatives, and structural review.

  \item[Human corrections]
    Corrections around evidence preservation, draft handling, state authority,
    LaTeX report value, unnecessary folder creation, deployment complexity, and
    missing browser validation.
\end{description}

\section{Status of this Document}

This document is intentionally human-facing. It should be used as a review
artifact, not as the source of future agent behavior. Future agents may read it
to understand why the architecture changed, but they must resolve operational
rules from stable files.

\begin{table}[htbp]
\centering
\caption{Authority status of related repository surfaces.}
\label{tab:authority-surfaces}
\small
\begin{tabularx}{\textwidth}{l X X}
\toprule
\textbf{Surface} & \textbf{Purpose} & \textbf{Authority Status} \\
\midrule
\safePath{agents/rules/core/} & Mission, constraints, public copy, execution modes, truth hierarchy & Authoritative \\
\safePath{agents/rules/workflows/} & Workflow contracts and allowed behavior per mode & Authoritative \\
\safePath{agents/rules/roles/} & Role boundaries and handoff rules & Authoritative \\
\safePath{agents/state/} & Current status, logs, summaries, research, evidence trail & Contextual evidence only \\
\safePath{agents/reports/} & Formal proof artifacts and human-readable retrospective reports & Evidence, not current rules \\
\safePath{agents/work/social/drafts/} & Candidate copy for human review & Candidate output only \\
\safePath{site/} and root entrypoints & Public runtime and static site surfaces & Product surface, not instruction source \\
\bottomrule
\end{tabularx}
\end{table}

\section{Non-Goals}

This report does not propose another large reorganization. It records the
architecture lessons needed before any such reorganization is safe. In
particular, it does not assume that fewer folders automatically means a better
system. The test is stricter: every path must have a unique reason to exist, and
no path should force manual synchronization of the same fact.

% ============================================================================
\chapter{Pre-Cleanup Failure Pattern}
\label{ch:pre-cleanup}
% ============================================================================

\takeaway{Before cleanup, the system had become too good at producing traces of
work and not reliable enough at preserving the user's requested scope. The
agentic harness was doing visible labor, but not always useful work.}

\section{Symptom 1: Documentation Loops}

The project treated documentation as a default closure obligation rather than a
tool used when it added value. Sprint plans, review files, closure matrices,
logs, pair-check notes, and public-narrative decisions could all be triggered
around a task that only needed a concrete content result.

The visible failure was an ``eternal documentation loop'': after the user
resolved blockers or declared human actions complete, the system still found new
process surfaces to close. The underlying issue was not documentation quality.
The issue was uncontrolled continuation.

\begin{description}[style=nextline]
  \item[Observed behavior]
    Agents kept producing run plans, closure notes, checklists, and validation
    traces even when the task was already complete.

  \item[Root cause]
    Workflow files and templates rewarded completion of process dimensions
    independent of whether those dimensions were relevant to the user's current
    request.

  \item[Correction]
    Tie documentation to execution mode. \workflow{text-only} and
    \workflow{answer-only} do not create sprint logs, state updates, or formal
    reports unless the user explicitly asks.
\end{description}

\section{Symptom 2: Over-Engineering Without Output}

The system often transformed small work into a full agentic event. That made the
project look active while delaying actual output. This mattered most during
public-copy and campaign work: a usable post, cover letter, or site paragraph
was often less important to the agent than satisfying the harness.

The warning sign is simple: when the harness can produce a convincing report
without producing the requested artifact, the harness is overfitted to its own
visibility.

\section{Symptom 3: Campaign Voice Drift}

The public voice drifted because several files competed to define it. The older
style-book material preserved useful history but also contained stale campaign
posture. Later site-voice corrections identified a better direction: less
sales-copy, less defensive symmetry, fewer generic ``not X but Y'' structures,
and more artifact-led, first-person, technical copy.

The problem was not that multiple voice artifacts existed. The problem was that
their authority status was unclear. If an old style guide and a current public
copy canon both exist without hierarchy, a model may choose the more verbose or
easier-to-follow file, even if it is outdated.

\section{Symptom 4: Tests and Code Changes for Text Work}

The system sometimes executed tests or modified implementation surfaces when the
user asked only for text. This is a classic missing-mode failure. A workflow
instruction like ``test where appropriate'' is safe only when ``appropriate'' is
defined by task type.

For this project, the expected boundary is:

\begin{itemize}
  \item text-only work does not run technical tests;
  \item site-copy work does not modify CSS, JavaScript, validators, or runtime
    behavior;
  \item implementation work validates the touched technical surface;
  \item system-refactor work validates rule consistency and affected
    infrastructure.
\end{itemize}

\section{Symptom 5: Manual Mirror Copies in the Site}

The site accumulated copies of repository files that already existed elsewhere:
status JSON, CV files, diagrams, screenshots, and assets. Some copies were exact
duplicates. Others had already diverged. This is worse than clutter: it creates a
false public surface.

If a user or agent updates the canonical CV but the site points to a stale copy,
the public site silently lies. The same applies to status files and report links.

\pitfall{Manual copy rules are not rules. They are future failures scheduled in
advance. If a file must be synchronized by memory, the architecture is wrong.}

% ============================================================================
\chapter{Cleanup Attempt and Second-Order Failures}
\label{ch:cleanup-failures}
% ============================================================================

\takeaway{The cleanup attempt revealed a second-order failure: a system designed
to reduce clutter can destroy evidence if it does not distinguish stale steering
from proof. The failure produced useful information because the human review
forced the missing categories into the architecture.}

\section{The First Correct Direction}

The initial diagnosis was materially correct. The repo needed:

\begin{itemize}
  \item execution modes;
  \item public-copy canon;
  \item read-only review by default;
  \item state as evidence/status, not authority;
  \item simplified PR and benchmark pressure;
  \item removal of stale style and tool-output sources;
  \item elimination of tracked site copies.
\end{itemize}

Those are still valid. The failure came from applying them too broadly and too
structurally.

\section{Error 1: Treating Formal Evidence as Clutter}

During cleanup, old reports and LaTeX artifacts were initially classified as
historical clutter. That was wrong. Formal reports are one of the public proof
surfaces of the project. They demonstrate documentation discipline, sprint
history, and professional communication. Removing them weakens the project.

The corrected distinction is:

\begin{description}[style=nextline]
  \item[Context poison]
    Superseded style guides, stale tool-output reports, discarded copy drafts,
    old instructions that compete with current rules.

  \item[Evidence]
    Formal reports, sprint logs, summaries, research, review notes, and state
    records that show what happened.

  \item[Current authority]
    Stable operating files under \safePath{agents/rules/}.
\end{description}

The correction was not ``keep everything.'' The correction was category-aware
retention.

\section{Error 2: Deleting Sprint Evidence}

A later cleanup removed sprint logs and summaries. This violated a central
project boundary: the repository is a proof loop. The evidence of past work must
remain available even when it is not used to steer future work.

The correct rule is precise:

\begin{quote}
State can be read for context and evidence. State cannot define future behavior,
voice, scope, acceptance, or strategy.
\end{quote}

This makes state safe without deleting it. If a past state entry is wrong, it is
wrong evidence; it is not a current rule.

\section{Error 3: Restoring Drafts as Evidence}

The next correction overcompensated by restoring old social drafts as evidence.
The user corrected this: discarded drafts are not evidence in the same way logs
are evidence. They are candidate outputs. Old discarded candidate text is a
high-risk source of copy poisoning because models may reuse its phrasing.

The final rule is:

\begin{itemize}
  \item future drafts may live in \safePath{agents/work/social/drafts/};
  \item drafts are human-readable candidate outputs only;
  \item drafts are not approved copy, blockers, strategy, or instruction;
  \item discarded old drafts should not remain as reusable context.
\end{itemize}

\section{Error 4: Adding Labels Instead of Fixing Authority}

The cleanup briefly tried to solve evidence safety by adding archive-marker files
and labels. That was unnecessary structure. The rule belongs in
\safePath{agents/rules/core/truth-hierarchy.md}, not in many small README markers under
state folders.

This revealed a broader principle: if a rule can be stated once in the
authoritative layer, do not distribute labels across evidence folders. Labels
are easy to forget, easy to duplicate, and easy for agents to over-read.

\section{Error 5: Solving Site Duplication with a New Generated Root}

The site-copy problem was real. The first solution, however, introduced a
generated \safePath{.site-dist/} folder at the repository root. That reduced
tracked duplication but added a new deployment concept and another path for
agents to understand.

The user challenge was valid: if the repository root already contains the
canonical sources, why generate another root-level copy? The later correction was
to simplify the deployment model rather than add an artifact layer.

\section{Error 6: Moving Too Much to Root}

The next correction moved more site entrypoint files to the root than the user
had clearly accepted. The user had said that having \safePath{index.html} in the
root was not so bad; that did not necessarily authorize moving every metadata
file or reshaping the whole public surface. This exposed a process failure, not
just a technical one.

The lesson is that questions must be treated as questions. When the user asks
``why this structure?'' or ``was this necessary?'', the next step is a design
review, not another implementation.

\pitfall{A cleanup run can fail by being too confident. When the task is
structural and the user is actively questioning assumptions, implementation
should pause at the tradeoff table.}

% ============================================================================
\chapter{Root Causes}
\label{ch:root-causes}
% ============================================================================

\takeaway{The repeated failures share one root: missing boundary algebra. The
system had files, roles, workflows, and evidence, but it did not always know
which of them could control the next action.}

\section{Authority Was Described, Not Enforced}

The repo already had a truth hierarchy, but some workflow and role files still
treated state as an operational input. That created a contradiction: stable docs
said one thing, but execution flows could still read state as a task contract,
integration checklist, blocker source, or acceptance source.

Authority must be enforced through verbs:

\begin{description}[style=nextline]
  \item[Docs define]
    rules, constraints, role boundaries, execution modes, workflow semantics.

  \item[State records]
    current status, evidence, logs, summaries, research, past decisions.

  \item[Content publishes]
    reports, social candidates, public pages, human-facing artifacts.

  \item[Site serves]
    public runtime, project pages, and canonical assets exposed through the repo.
\end{description}

\section{Execution Mode Was Missing}

Without execution modes, every task looked eligible for every workflow behavior.
This is why text work could trigger tests, site changes, logs, or sprint
closure. The system needs a compact mode table that every role respects.

\begin{table}[htbp]
\centering
\caption{Execution modes and default permissions.}
\label{tab:execution-modes}
\small
\begin{tabularx}{\textwidth}{l X X X}
\toprule
\textbf{Mode} & \textbf{Purpose} & \textbf{Allowed by Default} & \textbf{Forbidden by Default} \\
\midrule
\workflow{answer-only} & Respond without repository mutation & Explanation, recommendation, diagnosis & File edits, tests, state updates, commits \\
\workflow{text-only} & Produce or revise prose/copy & Declared text files or direct answer text & Code, CSS, JS, tests, generated data, state logs \\
\workflow{site-copy-only} & Change visible public copy inside existing site structure & Copy-bearing HTML/JSON/text files & Runtime JS, CSS, validators, deployment, assets unless requested \\
\workflow{implementation} & Change behavior, build, validation, or site runtime & Scoped code/config/docs changes & Unscoped rewrites, unrelated cleanup \\
\workflow{system-refactor} & Change stable rules or architecture & Explicitly declared stable files and validation plan & Opportunistic reorganization beyond the contract \\
\bottomrule
\end{tabularx}
\end{table}

\section{Review Was Too Easily Mutating}

Review should reduce risk. In this system, review could become another mutation
path. A read-only review that reports issues is a different operation from a
review that edits files, updates state, or triggers validators.

The safer rule is:

\begin{quote}
Review is read-only unless the user explicitly asks to fix, patch, implement, or
apply changes.
\end{quote}

This does not weaken review. It restores it as a decision tool.

\section{Benchmarks Rewarded Shape Instead of Behavior}

Some tests and benchmark definitions measured structural properties: counts,
sections, checklist shapes, or artifact presence. Those checks are useful only
when the structure is itself the product. In this case they created pressure to
keep ceremony.

Better tests would assert behavioral invariants:

\begin{itemize}
  \item text-only mode forbids code and test execution;
  \item state is never named as an authority source in stable docs;
  \item no tracked duplicate exists for canonical site-exposed assets;
  \item public-copy canon is referenced by social planning;
  \item old style-book and discarded draft paths are not live references.
\end{itemize}

\section{Validation Did Not Match Risk}

The project often ran broad validation because validation was available. That
does not mean it was relevant. Running Python tests after a copy edit adds noise.
Running link checks after a pure social post draft is unnecessary. Conversely,
site path changes require browser/runtime validation, not just static link
checks.

Validation should be selected by changed surface and execution mode, not by
habit.

\section{Structure Was Used as a Substitute for Design}

The cleanup attempts created or proposed new paths: active-run files, generated
site artifacts, archive labels, and root entrypoint moves. Some were reasonable
intermediate ideas. The issue was not that they were absurd. The issue was that
they were implemented or planned before the minimum viable architecture was
settled.

A structural change must pass three questions:

\begin{enumerate}
  \item Does this path hold information that cannot live cleanly in an existing
    path?
  \item Does it remove manual synchronization, or does it create another thing to
    keep aligned?
  \item Does it reduce future agent ambiguity, or does it add another route to
    reason about?
\end{enumerate}

If the answer is unclear, do not add the path yet.

% ============================================================================
\chapter{Corrected Architecture}
\label{ch:corrected-architecture}
% ============================================================================

\takeaway{The proper architecture is not a larger agentic harness. It is a small
set of hard boundaries: docs define rules, state records evidence, content holds
human-facing artifacts, drafts are candidate outputs, and the site exposes
canonical repository files without tracked duplication.}

\section{Minimal Repository Semantics}

The corrected architecture can be described without a new taxonomy explosion:

\begin{description}[style=nextline]
  \item[\safePath{AGENTS.md}]
    Entry point and routing index. It should stay compact. It should not become
    a second manual.

  \item[\safePath{agents/rules/core/}]
    Stable rules: mission, constraints, truth hierarchy, career direction,
    execution modes, public-copy canon, tool/CI policy when needed.

  \item[\safePath{agents/rules/workflows/}]
    Workflow behavior. Each workflow must respect execution mode and direct user
    scope.

  \item[\safePath{agents/rules/roles/}]
    Role boundaries. Roles should point to canonical rules instead of duplicating
    them.

  \item[\safePath{agents/rules/templates/}]
    Reusable shapes for real artifacts. Templates are not obligations.

  \item[\safePath{agents/state/}]
    Evidence, current status, logs, summaries, research, and backlog. Readable
    but non-authoritative.

  \item[\safePath{content/}]
    Human-facing content: formal reports, social drafts, research, and generated
    public proof.

  \item[\safePath{site/}]
    Site-owned runtime and content files only, unless the final structural
    review decides root-public files should live elsewhere.

  \item[\safePath{assets/} and \safePath{data/}]
    Canonical repository assets and machine-readable data when they are not
    intrinsically site-owned.
\end{description}

\section{Authority Flow}

\begin{figure}[htbp]
\centering
\begin{tikzpicture}[node distance=1.4cm and 2.0cm]
  \node[usernode] (user) {Direct user prompt};
  \node[filenode, below left=1.6cm and 1.6cm of user] (core) {agents/rules/core};
  \node[agentnode, below=of user] (workflows) {agents/rules/workflows};
  \node[agentnode, below right=1.6cm and 1.6cm of user] (roles) {agents/rules/roles};
  \node[statenode, below=2.4cm of workflows] (state) {state};
  \node[filenode, left=2.2cm of state] (content) {content/reports};
  \node[filenode, right=2.2cm of state] (site) {site / public surface};

  \draw[flowedge] (user) -- node[left,font=\scriptsize]{scope + mode} (core);
  \draw[flowedge] (user) -- node[right,font=\scriptsize]{task} (workflows);
  \draw[flowedge] (workflows) -- node[above,font=\scriptsize]{role contract} (roles);
  \draw[datadge] (core) -- node[left,font=\scriptsize]{rules} (workflows);
  \draw[datadge] (core) -- node[right,font=\scriptsize]{rules} (roles);
  \draw[datadge] (state) -- node[left,font=\scriptsize]{evidence only} (content);
  \draw[datadge] (state) -- node[right,font=\scriptsize]{status only} (site);
  \draw[datadge] (workflows) -- node[right,font=\scriptsize]{records} (state);
\end{tikzpicture}
\caption{Corrected authority flow: state can inform, but never command.}
\label{fig:authority-flow}
\end{figure}

\section{Drafts, Reports, and Evidence}

The corrected categories are intentionally narrow.

\begin{longtable}{p{0.22\textwidth} p{0.34\textwidth} p{0.34\textwidth}}
\caption{Retention policy by artifact type.}
\label{tab:retention-policy}\\
\toprule
\textbf{Artifact Type} & \textbf{Keep When} & \textbf{Remove or Demote When} \\
\midrule
\endfirsthead
\toprule
\textbf{Artifact Type} & \textbf{Keep When} & \textbf{Remove or Demote When} \\
\midrule
\endhead
Formal reports & They document completed work, decisions, or public proof & Only if duplicated, broken, or replaced by a newer formal artifact with explicit migration \\
Sprint logs and summaries & They record what happened and support auditability & Never use as current instructions; mark through hierarchy, not folder labels \\
Research reports & They contain market or technical evidence still useful for interpretation & Remove if purely superseded prompt output or stale strategy disguised as research \\
Social drafts & They are current candidate copy for human review & Delete discarded old drafts; never treat as approved or strategic \\
Style guides & They are the current canon under \safePath{agents/rules/core/} & Delete or demote superseded style files outside the rule layer \\
Tool-output reports & They contain final curated evidence & Quarantine raw tool output that steers future agents toward obsolete fixes \\
Site assets & They are site-owned runtime/design assets & Remove tracked mirror copies of canonical repo files \\
\bottomrule
\end{longtable}

\section{Site Source Model}

The site-source decision remained unsettled by the end of the conversation, and
that is important. The final implementation should not be inferred from one
intermediate correction. The stable conclusion is narrower:

\begin{itemize}
  \item canonical files should not be manually copied under \safePath{site/};
  \item if GitHub Pages serves the repository root, site paths may reference
    root canonical files directly;
  \item if GitHub Pages serves \safePath{site/}, the build/publish mechanism must
    expose required root files automatically;
  \item root-level public files should be minimal and justified;
  \item browser runtime validation is mandatory after site-path changes.
\end{itemize}

The final site topology should be decided by a separate design review, not by a
cleanup agent improvising during implementation.

\pitfall{The phrase ``just upload root'' does not automatically settle where
every site file belongs. It settles a deployment preference. File layout still
needs a small design decision.}

% ============================================================================
\chapter{Process Corrections for Future Runs}
\label{ch:process}
% ============================================================================

\takeaway{The next refactor should use less agentic ceremony, not more. The
process should force classification before mutation, preserve evidence before
deletion, and validate only the behavior that changed.}

\section{Decision Gate Before Any Structural Change}

When a user asks structural questions, the agent must not treat the question as
approval to edit. The next response should be a tradeoff table and a proposed
minimal action.

\begin{table}[htbp]
\centering
\caption{Required reasoning before structural edits.}
\label{tab:structure-gate}
\small
\begin{tabularx}{\textwidth}{l X X}
\toprule
\textbf{Question} & \textbf{Required Answer} & \textbf{Stop Condition} \\
\midrule
What problem does this path solve? & One specific maintenance, authority, or runtime issue & Stop if the answer is aesthetic only \\
Can an existing file hold the rule? & Identify the existing candidate and why it fails & Stop if existing file is sufficient \\
Does it remove duplication? & Show what manual sync disappears & Stop if it adds another sync target \\
Will agents understand it faster? & Explain reduced routing/context load & Stop if it creates another route \\
Can it be validated? & Name a concrete check & Stop if validation is vague \\
\bottomrule
\end{tabularx}
\end{table}

\section{File Triage Before Cleanup}

Cleanup must classify every candidate before deletion.

\begin{description}[style=nextline]
  \item[Delete]
    Discarded drafts, superseded active-style rules, raw tool outputs that
    mislead future agents, duplicate tracked site copies, generated caches.

  \item[Keep]
    Formal reports, logs, summaries, research that records decisions, current
    public site/runtime files, canonical assets, current data.

  \item[Demote]
    Historical references that are useful to humans but not current rules. The
    demotion must be encoded in the truth hierarchy, not with many local labels.

  \item[Review]
    Any file that is both evidence and possible poison. Do not delete until its
    evidence is preserved elsewhere or its role is explicitly resolved.
\end{description}

\section{Validation by Mode}

\begin{longtable}{p{0.22\textwidth} p{0.34\textwidth} p{0.34\textwidth}}
\caption{Validation plan by touched surface.}
\label{tab:validation-by-mode}\\
\toprule
\textbf{Touched Surface} & \textbf{Required Validation} & \textbf{Not Required by Default} \\
\midrule
\endfirsthead
\toprule
\textbf{Touched Surface} & \textbf{Required Validation} & \textbf{Not Required by Default} \\
\midrule
\endhead
Direct answer only & None beyond reasoning check & Repo tests, link checks, state updates \\
Social/copy draft & Human readability and canon check & Pytest, static-site validator, browser check \\
Markdown docs & Link/reference check if paths changed & Browser tests unless site runtime changed \\
Stable rule files & Rule contradiction scan, targeted benchmarks & Full site render unless public runtime changed \\
Site content paths & Static validator and browser smoke check & LaTeX build unless report source changed \\
Site runtime JS/CSS & Browser smoke check, static validator & Unrelated Python tests unless scripts changed \\
Python scripts/tests & Targeted pytest & LaTeX/browser unless affected \\
LaTeX source & LaTeX build for touched target & Site/browser unless links changed public surface \\
\bottomrule
\end{longtable}

\section{Browser Validation Is a Separate Class}

The conversation exposed a serious validation gap: many checks can pass while
the website is broken. Static path validation is not browser validation. If
routing, root placement, module paths, fetch paths, or public entrypoints change,
the run must load the site through a local HTTP server and inspect:

\begin{itemize}
  \item page title and main route render;
  \item JavaScript console errors;
  \item failed network requests;
  \item representative internal navigation;
  \item responsive sanity if layout changed.
\end{itemize}

The validation does not need to be theatrical. It needs to test the runtime.

\section{Evidence Preservation Rule}

Before deleting any historical file, ask:

\begin{enumerate}
  \item Is this an output candidate, an instruction source, or evidence?
  \item If evidence, is it already represented in a formal report?
  \item If not, does deletion remove proof of work?
  \item If it is poisonous, can it be made non-authoritative by hierarchy instead
    of deletion?
  \item If it is a discarded draft, is there any reason to preserve the body?
\end{enumerate}

Only discarded drafts, duplicate generated files, stale active instructions, and
raw misleading tool output should be easy deletes. Evidence requires a higher
bar.

% ============================================================================
\chapter{Recommended Minimal Refactor Plan}
\label{ch:recommended-plan}
% ============================================================================

\takeaway{The next successful run should be boring: classify, patch only
authoritative files, remove only proven duplication, browser-test site changes,
and record the result once.}

\section{Phase 1: Freeze the Authority Model}

\begin{enumerate}
  \item Confirm \safePath{agents/rules/core/truth-hierarchy.md} states that
    \safePath{agents/state/} is evidence/status only.
  \item Confirm \safePath{agents/rules/core/execution-modes.md} exists or add it as the
    only new core rule file.
  \item Confirm \safePath{agents/rules/core/public-copy.md} is the only current voice
    canon.
  \item Remove or demote any stable-doc wording that makes
    \safePath{agents/state/active-sprint.md} a contract or acceptance source.
  \item Replace duplicated bounded-contract fields in role files with references
    to \safePath{agents/rules/core/constraints.md}.
\end{enumerate}

\section{Phase 2: Collapse Process Ceremony}

\begin{enumerate}
  \item Make \workflow{review} read-only by default.
  \item Require explicit user wording before mutating review, state writes, or
    cleanup edits.
  \item Simplify PR template checks by execution mode.
  \item Rewrite or quarantine benchmarks that count checklist sections instead
    of behavior.
  \item Remove any rule that forces public-narrative, formal documentation, or
    backlog output for tasks where those are not relevant.
\end{enumerate}

\section{Phase 3: Clean Source Duplication Without Moving the World}

\begin{enumerate}
  \item List every tracked duplicate and classify whether it is exact,
    divergent, or merely similarly named.
  \item Delete tracked mirror copies only after all live references point to the
    canonical file.
  \item Do not delete unused assets merely because they are unused, unless they
    are mirror copies, stale public outputs, or obvious generated trash.
  \item For site architecture, decide root-vs-site serving with a design note
    before moving files.
  \item If root serving is selected, keep root entrypoints minimal and document
    exactly why each root public file exists.
\end{enumerate}

\section{Phase 4: Validate What Actually Changed}

The minimum validation plan for a structural/site refactor is:

\begin{itemize}
  \item path/reference scan for deleted or moved files;
  \item static-site validator;
  \item browser load through local HTTP server;
  \item targeted pytest only if Python/tests/benchmarks changed;
  \item LaTeX build only if report sources changed;
  \item no broad validation for pure text-only work.
\end{itemize}

\section{Phase 5: Record Once}

The output of the run should be one evidence note, not another forest of logs.
That note should record:

\begin{itemize}
  \item changed authority rules;
  \item deleted files and reason categories;
  \item preserved evidence categories;
  \item validation performed;
  \item unresolved design questions.
\end{itemize}

It should not become a new instruction source.

% ============================================================================
\chapter{Open Structural Questions}
\label{ch:open-questions}
% ============================================================================

\takeaway{Some structural questions remain legitimate. They should be answered
through a small design decision, not through another cleanup implementation.}

\section{Should the Site Live at Root?}

The user raised the concern that moving more than \safePath{index.html} to root
may add clutter. This is a valid architectural question.

\begin{table}[htbp]
\centering
\caption{Site-root alternatives.}
\label{tab:site-root-options}
\small
\begin{tabularx}{\textwidth}{l X X}
\toprule
\textbf{Option} & \textbf{Benefit} & \textbf{Risk} \\
\midrule
Serve \safePath{site/} & Keeps public runtime isolated from repo root & Requires a publish mechanism to expose canonical root files \\
Serve repo root & Lets browser access canonical \safePath{assets/}, \safePath{data/}, and reports directly & Mixes public entrypoints with repo operating files \\
Root \safePath{index.html} only & Smallest visible move & May still require path redirects or loader indirection \\
Move full site to root & Simplest static hosting mental model & Pollutes root and collides with repo-as-system organization \\
\bottomrule
\end{tabularx}
\end{table}

No option is free. The current recommendation is to choose the model that removes
manual duplication with the fewest new concepts, then validate in a real browser.

\section{Why Separate \safePath{assets/}, \safePath{content/}, and \safePath{data/}?}

The distinction can be justified, but only if it remains clear:

\begin{description}[style=nextline]
  \item[\safePath{assets/}]
    Binary or design-support materials: images, CV PDFs, diagrams, screenshots.
    These are files, not structured content.

  \item[\safePath{content/}]
    Human-facing authored artifacts: reports, social drafts, research, long-form
    publication material.

  \item[\safePath{data/}]
    Machine-readable facts used by tooling or site runtime, such as public
    status JSON.
\end{description}

If the repo cannot explain that distinction in one paragraph, it is too
complicated. If files regularly migrate between these categories, the categories
need revision.

\section{Why Rules Now Live Under \safePath{agents/}}

The implemented model keeps all agent-read and agent-produced material under
\safePath{agents/}. Stable rules live in \safePath{agents/rules/}; evidence and
status live in \safePath{agents/state/}; report sources live in
\safePath{agents/reports/tex/}; tools and tests live in \safePath{agents/tools/}
and \safePath{agents/tests/}. This makes the root smaller and prevents public
runtime files from being mixed with internal operating material.

The authority boundary is not ``agents can do whatever is under
\safePath{agents/}.'' It is stricter: \safePath{agents/rules/} is authoritative,
while \safePath{agents/state/} is evidence only.

\section{What Should Happen to Unused Assets?}

Unused is not the same as harmful. The cleanup should prioritize duplicated
source-of-truth risk over unused future material. Unused generated images or logo
variants can remain if they do not create manual synchronization rules or
mislead agents. They are low priority.

\pitfall{Minimalism is not deletion. Minimalism is fewer active meanings. A file
can exist as evidence or optional material without becoming a rule.}

% ============================================================================
\chapter{Acceptance Criteria for the Next Correct Run}
\label{ch:acceptance}
% ============================================================================

\takeaway{A successful follow-up run is not the one with the largest diff. It is
the one that reduces ambiguity, preserves evidence, removes synchronization
risk, and validates the actual runtime if touched.}

\section{Required Outcomes}

\begin{enumerate}
  \item No stable file describes \safePath{agents/state/} as authoritative for rules,
    acceptance, behavior, or voice.
  \item Text-only and site-copy-only modes have explicit negative write scopes.
  \item Review is read-only by default.
  \item Current public-copy canon lives in \safePath{agents/rules/core/}.
  \item Superseded style and discarded old draft bodies are not live references.
  \item Formal reports and state logs remain available as evidence.
  \item No tracked manual mirror copies exist for site-exposed canonical assets
    unless a specific exception is documented.
  \item Site-path changes are validated in a real browser, not only by link
    scans.
  \item Any new folder has a stated necessity test and cannot be replaced by an
    existing path.
  \item The run closes with one evidence record and does not create a new
    documentation loop.
\end{enumerate}

\section{Failure Conditions}

The run should be considered failed or partial if it:

\begin{itemize}
  \item deletes evidence without preserving or classifying it;
  \item creates new paths only to satisfy the harness;
  \item runs broad tests for text-only changes;
  \item changes site runtime without browser validation;
  \item reintroduces manual copy requirements;
  \item treats user questions as implementation approval;
  \item leaves stale references to deleted drafts or superseded style files;
  \item adds local labels where a single core rule would suffice.
\end{itemize}

\section{Operational Checklist}

\begin{longtable}{p{0.06\textwidth} p{0.40\textwidth} p{0.44\textwidth}}
\caption{Minimal follow-up checklist.}
\label{tab:followup-checklist}\\
\toprule
\textbf{\#} & \textbf{Check} & \textbf{Evidence} \\
\midrule
\endfirsthead
\toprule
\textbf{\#} & \textbf{Check} & \textbf{Evidence} \\
\midrule
\endhead
1 & Classify execution mode before action & Stated mode and allowed writes \\
2 & Classify every deletion candidate & Delete/keep/demote/review table \\
3 & Patch authority in stable docs only & Diff limited to necessary rule files \\
4 & Remove manual duplicates after reference update & Reference scan and duplicate hash check \\
5 & Validate by changed surface & Mode-specific validation log \\
6 & Browser-test if runtime changed & Console/network/render evidence \\
7 & Preserve formal reports and state evidence & No deleted evidence paths unless intentionally migrated \\
8 & Close with one evidence note & Single retrospective or run note \\
\bottomrule
\end{longtable}

% ============================================================================
\chapter{Conclusion}
\label{ch:conclusion}
% ============================================================================

\takeaway{The cleanup failure is useful because it exposed the real work:
designing an agentic system that can simplify itself without erasing its own
proof, expanding its own process, or mistaking confidence for correctness.}

AgenticCareerBoost remains valuable precisely because these failures are visible.
A polished portfolio would hide this class of engineering problem. This
repository exposes it: agents can be useful, but only when authority, scope,
evidence, and validation are explicit enough to resist drift.

The final architecture should not be a grander agentic operating system. It
should be a smaller one with harder edges:

\begin{itemize}
  \item direct user scope controls execution mode;
  \item docs define rules;
  \item state records evidence;
  \item content preserves human-facing proof;
  \item drafts remain candidate outputs;
  \item the site shows canonical repository material;
  \item validation follows the changed surface;
  \item structural changes require reasoning before implementation.
\end{itemize}

The next successful refactor should therefore feel less impressive while it is
running. Fewer agents, fewer logs, fewer invented paths, fewer validations, and
more exact reasoning. That is the standard this report records.

\pitfall{Do not optimize the repository to look agentic. Optimize it so agents
have fewer ways to misunderstand what matters.}

\appendix

% ============================================================================
\chapter{Compact Decision Ledger}
\label{app:decision-ledger}
% ============================================================================

\begin{longtable}{p{0.18\textwidth} p{0.34\textwidth} p{0.38\textwidth}}
\caption{Decisions extracted from the refactor discussion.}
\label{tab:decision-ledger}\\
\toprule
\textbf{Decision} & \textbf{Accepted Direction} & \textbf{Reason} \\
\midrule
\endfirsthead
\toprule
\textbf{Decision} & \textbf{Accepted Direction} & \textbf{Reason} \\
\midrule
\endhead
Authority & \safePath{agents/rules/} is authoritative; \safePath{agents/state/} is evidence/status & Prevents poisoned historical state from steering future runs \\
Evidence & Keep logs, summaries, reports, and research & They prove work happened and support auditability \\
Drafts & Keep future draft folder, delete discarded old draft bodies & Drafts are candidate outputs, not approved evidence or rules \\
Voice & Canonical public copy belongs in \safePath{agents/rules/core/public-copy.md} & Prevents stale style-book material from overriding current direction \\
Review & Read-only by default & Prevents review from becoming another mutation path \\
Validation & Mode-specific & Avoids broad tests for text-only work and missing browser tests for site work \\
Site assets & Do not track manual mirror copies & Prevents stale public files and synchronization burden \\
New folders & Require necessity test & Prevents cleanup from adding more clutter than it removes \\
Root site & Open design question & Root upload may simplify serving but can pollute repo root if over-applied \\
\bottomrule
\end{longtable}

% ============================================================================
\chapter{Suggested Placement}
\label{app:placement}
% ============================================================================

This source is intended to live under:

\begin{center}
\safePath{agents/reports/tex/sprints/agentic-system-refactor-retrospective.tex}
\end{center}

A future PDF build can be promoted to:

\begin{center}
\safePath{site/files/reports/agentic-system-refactor-retrospective.pdf}
\end{center}

If included in repository status, it should be described as evidence, not as a
new canonical rule source.

\end{document}
